% ============================================================
% INFORME DEL COMPILADOR HULK
% Sección 1: Portada + Resumen/Abstract
% ============================================================

\documentclass[12pt, a4paper]{article}

% --- Paquetes esenciales ---
\usepackage[utf8]{inputenc}
\usepackage[T1]{fontenc}
\usepackage[spanish]{babel}
\usepackage{geometry}
\geometry{margin=2.5cm}
\usepackage{hyperref}
\hypersetup{
    colorlinks=true,
    linkcolor=blue,
    urlcolor=blue,
    citecolor=blue
}
\usepackage{graphicx}
\usepackage{amsmath}
\usepackage{amssymb}
\usepackage{listings}
\usepackage{xcolor}
\usepackage{tikz}
\usetikzlibrary{arrows.meta, positioning}
\usepackage{booktabs}
\usepackage{array}
\usepackage{tabularx}
\usepackage{enumitem}
\usepackage{setspace}
\usepackage{fancyhdr}
\usepackage{titlesec}
\usepackage{caption}
\usepackage{subcaption}
\usepackage{cite}
\usepackage{url}

% --- Configuración de listings para código HULK ---
\definecolor{hulkgreen}{rgb}{0.0, 0.5, 0.0}
\definecolor{hulkgray}{rgb}{0.5, 0.5, 0.5}
\definecolor{hulkblue}{rgb}{0.0, 0.0, 0.7}
\definecolor{hulkpurple}{rgb}{0.5, 0.0, 0.5}
\definecolor{hulkbg}{rgb}{0.97, 0.97, 0.97}

\lstdefinelanguage{HULK}{
    keywords={function, type, protocol, inherits, let, in, if, elif, else,
              while, for, new, self, base, true, false, is, as, match, case},
    keywordstyle=\color{hulkblue}\bfseries,
    comment=[l]{//},
    commentstyle=\color{hulkgray}\itshape,
    stringstyle=\color{hulkgreen},
    morestring=[b]",
    sensitive=true
}

\lstset{
    language=HULK,
    basicstyle=\ttfamily\small,
    backgroundcolor=\color{hulkbg},
    frame=single,
    rulecolor=\color{hulkgray},
    numbers=left,
    numberstyle=\tiny\color{hulkgray},
    numbersep=5pt,
    breaklines=true,
    showstringspaces=false,
    tabsize=4,
    captionpos=b,
    literate=
        {á}{{\'{a}}}1 {é}{{\'{e}}}1 {í}{{\'{i}}}1
        {ó}{{\'{o}}}1 {ú}{{\'{u}}}1 {ñ}{{\~{n}}}1
        {Á}{{\'{A}}}1 {É}{{\'{E}}}1 {Í}{{\'{I}}}1
        {Ó}{{\'{O}}}1 {Ú}{{\'{U}}}1 {Ñ}{{\~{N}}}1
}

% --- Encabezado y pie de página ---
\pagestyle{fancy}
\fancyhf{}
\fancyhead[L]{\small Compilador HULK}
\fancyhead[R]{\small Universidad de La Habana, 2026}
\fancyfoot[C]{\thepage}

% --- Título personalizado ---
\titleformat{\section}
  {\normalfont\Large\bfseries\color{hulkblue}}
  {\thesection}{1em}{}

\titleformat{\subsection}
  {\normalfont\large\bfseries}
  {\thesubsection}{1em}{}

\titleformat{\subsubsection}
  {\normalfont\normalsize\bfseries}
  {\thesubsubsection}{1em}{}

\lstdefinelanguage{Rust}{
    keywords={pub, fn, let, mut, struct, enum, impl, use, mod, match,
              if, else, for, while, return, true, false, self, super,
              type, trait, where, async, await, move, ref, in, loop},
    keywordstyle=\color{hulkpurple}\bfseries,
    comment=[l]{//},
    commentstyle=\color{hulkgray}\itshape,
    stringstyle=\color{hulkgreen},
    morestring=[b]",
    sensitive=true
}
% ============================================================
\begin{document}

% ============================================================
% PORTADA
% ============================================================
\begin{titlepage}
    \centering
    \vspace*{1cm}

    {\Huge\bfseries\color{hulkblue}
        Compilador del Lenguaje HULK\\[0.5em]
        \Large Diseño, Implementación y Extensiones
    }

    \vspace{1.5cm}

    \rule{\linewidth}{0.5mm}

    \vspace{1cm}

    {\large
        \textbf{Informe de Proyecto}\\[0.3em]
        Asignaturas: \textit{Lenguajes de Programación} y \textit{Compilación}
    }

    \vspace{1.5cm}

    {\large\bfseries Autores:}\\[0.5em]
    \begin{tabular}{ll}
        Darío Francisco Alfonso Urrutia   & \texttt{@D4R102004}  \\[0.3em]
        Juan Carlos Carmenate Díaz       & \texttt{@Juank404}   \\[0.3em]
        Sebastian González Alfonso       & \texttt{@sebagonz106}
    \end{tabular}

    \vspace{1.5cm}

    {\large
        Facultad de Matemática y Computación\\
        Universidad de La Habana\\[0.5em]
        Junio de 2026
    }

    \vspace{1.5cm}

    \rule{\linewidth}{0.5mm}

    \vspace{1cm}

    {\small
        Repositorio: \url{https://github.com/D4R102004/hulk-compiler}
    }

\end{titlepage}

% ============================================================
% RESUMEN Y ABSTRACT
% ============================================================
\thispagestyle{empty}
\vspace*{1cm}

\begin{center}
    {\Large\bfseries Resumen / Abstract}
\end{center}

\vspace{1em}

\noindent\textbf{Resumen}

\vspace{0.5em}

\noindent
Este informe describe el diseño e implementación de un compilador completo para el
lenguaje HULK (\textit{Havana University Language for Kompilers}), desarrollado como
proyecto final de las asignaturas Lenguajes de Programación y Compilación de la
Universidad de La Habana. El compilador está implementado íntegramente en Rust e incluye
un analizador léxico, un analizador sintáctico descendente recursivo LL(1), un
analizador semántico de cinco pasadas con inferencia de tipos, y un generador de
código basado en LLVM 17 mediante la biblioteca \texttt{inkwell}. El compilador cubre
todos los aspectos del lenguaje HULK hasta la verificación de tipos (sección 16.8 de
la documentación oficial) e introduce tres extensiones al lenguaje: \textit{expresiones
\textbf{match}} para coincidencia de patrones estructurada, \textit{tipos función} que
permiten tratar funciones y métodos como valores de primera clase, y \textit{vectores}
con soporte para comprensiones al estilo funcional. El informe analiza las decisiones
de diseño más relevantes, las compara con soluciones adoptadas en lenguajes como
Kotlin, Rust, Haskell y Python, y documenta las pruebas de validación realizadas.

\vspace{1.5em}

\noindent\textbf{Abstract}

\vspace{0.5em}

\noindent
This report describes the design and implementation of a complete compiler for the
HULK language (\textit{Havana University Language for Kompilers}), developed as the
final project for the Programming Languages and Compilation courses at the University
of Havana. The compiler is written entirely in Rust and comprises a hand-written lexer,
a hand-written LL(1) recursive-descent parser, a five-pass semantic analyzer with
type inference, and a code generator targeting LLVM 17 via the \texttt{inkwell} crate.
The compiler covers all HULK features through type verification (§16.8) and introduces
three language extensions: \textbf{match expressions} for structured pattern matching,
\textbf{function types} enabling first-class function values, and \textbf{vectors}
with functional-style comprehension syntax. The report analyzes the key design
decisions, compares them to solutions in Kotlin, Rust, Haskell, and Python, and
documents the validation test suite.

\vspace{2em}

\noindent\textbf{Palabras clave:} compiladores, HULK, Rust, LLVM, inferencia de tipos,
coincidencia de patrones, funciones de primera clase, vectores.

\newpage

% ============================================================
% TABLA DE CONTENIDOS
% ============================================================
\tableofcontents
\newpage


% ============================================================
% SECCIÓN 2: INTRODUCCIÓN
% ============================================================

\section{Introducción}
\label{sec:intro}

\subsection{Historia de la compilación y los lenguajes de programación}

El nacimiento de los compiladores está íntimamente ligado al origen de la
programación moderna. En 1952, Grace Hopper desarrolló el primer compilador
conocido, el A-0 System, que traducía código simbólico matemático a instrucciones
de máquina. Este hito marcó el inicio de una nueva era: por primera vez, los
programadores podían expresar algoritmos en un nivel de abstracción superior al
del lenguaje ensamblador.

El salto definitivo llegó en 1957 con FORTRAN (\textit{Formula Translation}),
desarrollado por John Backus y su equipo en IBM. FORTRAN demostró que era posible
generar código máquina tan eficiente como el escrito a mano, disipando las dudas
de quienes creían que la abstracción debía sacrificar rendimiento. A FORTRAN le
siguió LISP en 1958, introduciendo la recursión y el tratamiento de funciones como
objetos de primera clase ---conceptos que siguen siendo relevantes en el diseño de
lenguajes actuales.

Las décadas de 1960 y 1970 vieron la consolidación de la teoría formal detrás de
la compilación. Noam Chomsky formuló la jerarquía de gramáticas formales, mientras
que Donald Knuth definió los lenguajes LR(k), sentando las bases matemáticas del
análisis sintáctico. La publicación del ``libro del dragón''
(\textit{Compilers: Principles, Techniques, and Tools}) por Aho, Sethi y Ullman
en 1986 codificó el conocimiento acumulado en una referencia canónica que aún se
utiliza en cursos universitarios de todo el mundo.

El estado del arte moderno incluye compiladores de infraestructura compleja como
GCC y LLVM. Este último, desarrollado originalmente por Chris Lattner en 2003,
introdujo un concepto transformador: una \textit{representación intermedia} (IR)
de bajo nivel y tipada que permite separar completamente el \textit{front end}
(análisis y tipado) del \textit{back end} (generación de código nativo). Hoy en
día, LLVM es la columna vertebral de compiladores para lenguajes como Clang
(C/C++), Swift, Kotlin Native y Rust.

\subsection{El lenguaje HULK y el contexto del proyecto}

HULK (\textit{Havana University Language for Kompilers}) es un lenguaje de
programación didáctico diseñado en la Facultad de Matemática y Computación de la
Universidad de La Habana, concebido específicamente para el estudio de la
compilación. HULK es un lenguaje de tipado estático, orientado a objetos y basado
en expresiones: toda construcción sintáctica ---incluyendo bloques, condicionales y
bucles--- produce un valor.

El lenguaje incluye un sistema de tipos con inferencia, herencia simple,
polimorfismo, protocolos (equivalentes a interfaces estructurales), iterables y
vectores. Su diseño deliberadamente compacto permite a un equipo de estudiantes
implementar un compilador completo en el transcurso de un semestre, cubriendo
todas las fases clásicas: análisis léxico, sintáctico, semántico y generación de
código.

El presente proyecto, desarrollado como trabajo final de las asignaturas
Lenguajes de Programación y Compilación, consiste en la implementación de un
compilador completo para HULK en el lenguaje Rust, generando código nativo para
Linux x86\_64 mediante el \textit{backend} LLVM 17.

\subsection{Motivación y objetivos de las extensiones}

La especificación base de HULK cubre los mecanismos fundamentales de la
programación orientada a objetos con tipos estáticos. Sin embargo, el lenguaje
original carece de herramientas expresivas modernas que permiten escribir código
más conciso, seguro y composable. Este proyecto introduce tres extensiones al
lenguaje, cada una motivada por necesidades concretas:

\begin{enumerate}[leftmargin=*, label=\arabic*.]
    \item \textbf{Expresiones \texttt{match}.}
    La ausencia de coincidencia de patrones obliga a escribir cadenas de
    \texttt{if}/\texttt{elif}/\texttt{else} verbosas y propensas a errores.
    La expresión \texttt{match} permite inspeccionar el valor y el tipo de una
    expresión de forma declarativa, siguiendo el modelo de lenguajes como Rust,
    Kotlin y Haskell.

    \item \textbf{Tipos función y funciones de primera clase.}
    En HULK base, las funciones son entidades globales y no pueden pasarse como
    argumentos ni almacenarse en variables. Esta extensión introduce el tipo
    \texttt{(T$_1$, ..., T$_n$) -> R} para representar funciones como valores,
    habilitando la programación funcional de alto orden: paso de funciones como
    argumentos, devolución de funciones desde funciones y referencias a métodos.

    \item \textbf{Vectores con comprensiones.}
    Las colecciones con tipo homogéneo son esenciales en cualquier lenguaje
    práctico. Esta extensión añade el tipo \texttt{Vector<T>} con sintaxis de
    literal (\texttt{[1, 2, 3]}) e indexación, así como comprensiones de vectores
    al estilo funcional (\texttt{[expr | var in iterable]}).
\end{enumerate}

Los tres objetivos de diseño que guían estas extensiones son: (1)
\textit{consistencia} con la semántica basada en expresiones de HULK,
(2) \textit{seguridad de tipos} verificada estáticamente, y (3)
\textit{expresividad} comparable a la de lenguajes modernos como Kotlin y Rust.

\subsection{Estructura del informe}

El resto del informe se organiza de la siguiente manera. La
Sección~\ref{sec:hulk} describe el lenguaje HULK original, sus paradigmas y su
sistema de tipos. La Sección~\ref{sec:extensiones} presenta en detalle las tres
extensiones propuestas, con su sintaxis formal, reglas de tipado, semántica
dinámica y comparativas con otros lenguajes. La Sección~\ref{sec:arquitectura}
analiza la arquitectura completa del compilador y justifica las decisiones de
diseño más relevantes. La Sección~\ref{sec:pruebas} documenta la metodología de
pruebas y los casos más significativos. La Sección~\ref{sec:discusion} discute
las dificultades encontradas y las alternativas descartadas. Finalmente, la
Sección~\ref{sec:conclusiones} presenta las conclusiones y el trabajo futuro.

\newpage

% ============================================================
% SECCIÓN 3: DESCRIPCIÓN GENERAL DEL LENGUAJE HULK
% Pegar antes de \end{document}
% ============================================================

\section{Descripción general del lenguaje HULK}
\label{sec:hulk}

\subsection{Paradigma y filosofía del lenguaje}

HULK (\textit{Havana University Language for Kompilers}) es un lenguaje de
programación estáticamente tipado, orientado a objetos y \textbf{basado en
expresiones}. A diferencia de lenguajes como C o Java, donde las construcciones
de control de flujo son \textit{sentencias} (sin valor de retorno), en HULK toda
construcción produce un valor: los bloques, los condicionales, los bucles y la
coincidencia de patrones son expresiones que pueden aparecer en cualquier posición
donde se espere un valor. Esta filosofía, compartida con lenguajes como Rust,
Scala y Kotlin, simplifica la semántica del lenguaje y fomenta un estilo de
programación más declarativo.

El programa HULK más sencillo ilustra esta idea:

\begin{lstlisting}[caption={Programa mínimo en HULK}, label=lst:hello]
print("Hello, World!");
\end{lstlisting}

Un programa HULK está compuesto por una secuencia de \textit{declaraciones
globales} (funciones, tipos y protocolos) seguida de una única
\textit{expresión de entrada} que constituye el punto de inicio de la ejecución.

\subsection{Sistema de tipos}

HULK implementa un sistema de tipos estático con inferencia parcial. Los tipos
se organizan en una jerarquía nominal con \texttt{Object} como raíz, y el
sistema admite tanto anotaciones explícitas como inferencia automática en la
mayoría de los contextos.

\subsubsection{Tipos primitivos y jerarquía}

El sistema de tipos del compilador implementado define los siguientes tipos
fundamentales, representados como variantes del enum \texttt{Type} en el módulo
\texttt{hulk-semantic}:

\begin{table}[h!]
\centering
\caption{Tipos primitivos del sistema de tipos de HULK}
\label{tab:types}
\begin{tabular}{@{}lll@{}}
\toprule
\textbf{Tipo HULK} & \textbf{Variante interna} & \textbf{Descripción} \\
\midrule
\texttt{Number}   & \texttt{Type::Number}   & Número de punto flotante (f64) \\
\texttt{String}   & \texttt{Type::String}   & Cadena de texto inmutable \\
\texttt{Boolean}  & \texttt{Type::Boolean}  & Valor lógico (\texttt{true}/\texttt{false}) \\
\texttt{Object}   & \texttt{Type::Object}   & Raíz de la jerarquía nominal \\
\texttt{T} (usuario) & \texttt{Type::Named(String)} & Tipo o protocolo definido por el usuario \\
\texttt{Vector<T>} & \texttt{Type::Vector(Box<Type>)} & Colección homogénea (extensión) \\
\texttt{T*}       & \texttt{Type::Iterable(Box<Type>)} & Protocolo iterable \\
\texttt{(T)->R}   & \texttt{Type::Function\{...\}} & Tipo función (extensión) \\
\bottomrule
\end{tabular}
\end{table}

Los tipos \texttt{Number}, \texttt{String} y \texttt{Boolean} son tipos de
valor primitivos: no pueden ser heredados y se representan sin boxing en la
generación de código. \texttt{Object} es el supertipo universal — todo tipo
HULK conforma a \texttt{Object}. Los tipos \texttt{Vector<T>} y
\texttt{(T)->R} son extensiones introducidas en este proyecto, descritas en
detalle en la Sección~\ref{sec:extensiones}.

La relación de conformidad (\textit{subtyping}) $T_1 \leq T_2$ se implementa
en el método \texttt{conforms\_to} y sigue las reglas:

\begin{enumerate}[leftmargin=*, label=\arabic*.]
    \item \textbf{Reflexividad:} $T \leq T$ para todo $T$.
    \item \textbf{Top:} $T \leq \texttt{Object}$ para todo $T$.
    \item \textbf{Herencia nominal:} $T \leq U$ si $U$ es ancestro de $T$ en la
          jerarquía de herencia.
    \item \textbf{Conformidad estructural a protocolos:} $T \leq P$ (protocolo)
          si $T$ implementa todos los métodos de $P$ con las signaturas correctas.
    \item \textbf{Covarianza de \texttt{Vector}:} $\texttt{Vector}<T> \leq
          \texttt{Vector}<U>$ si $T \leq U$.
    \item \textbf{Tipos función:} $(A_1,...,A_n) \to R_1 \leq (B_1,...,B_n) \to R_2$
          si $B_i \leq A_i$ (contravarianza en parámetros) y $R_1 \leq R_2$
          (covarianza en retorno).
\end{enumerate}

\subsubsection{Inferencia de tipos}

HULK permite declarar funciones sin anotaciones de tipo explícitas. El
compilador infiere los tipos a partir del uso:

\begin{lstlisting}[caption={Inferencia de tipos en HULK}, label=lst:inference]
// Sin anotaciones: el compilador infiere Number -> Number
function square(x) {
    x * x;
}

// Con anotaciones explícitas
function add(x: Number, y: Number): Number {
    x + y;
}
\end{lstlisting}

El análisis semántico implementado resuelve la inferencia de tipos en cinco
pasadas ordenadas (ver Sección~\ref{sec:arquitectura}).

\subsection{Características principales del lenguaje}

\subsubsection{Variables y ámbito léxico}

HULK usa \texttt{let ... in} como construcción de enlace de variables. El
ámbito es léxico y las variables pueden ser shadowed en ámbitos internos:

\begin{lstlisting}[caption={Variables y shadowing en HULK}, label=lst:let]
let x = 10 in {
    if (x == 10) print("ok") else print("fail");
    let x = 20 in {       // x queda en la sombra (shadowing)
        if (x == 20) print("ok") else print("fail");
    };
    if (x == 10) print("ok") else print("fail");
};
\end{lstlisting}

La asignación destructiva \texttt{:=} modifica el valor de una variable
existente sin crear un nuevo enlace:

\begin{lstlisting}[caption={Asignación destructiva}, label=lst:assign]
let i = 0 in
let result = 0 in {
    while (i < 5) {
        result := result + i;
        i := i + 1;
    };
    print(result);   // imprime 10
};
\end{lstlisting}

\subsubsection{Funciones}

Las funciones en HULK son globales y admiten dos formas sintácticas: la forma
\textit{inline} con \texttt{=>} y la forma de bloque con \texttt{\{\}}. Soportan
recursión y referencias mutuas (gracias al paso de recolección de declaraciones
que registra todas las signaturas antes de inferir los cuerpos):

\begin{lstlisting}[caption={Recursión mutua en HULK}, label=lst:mutual]
function is_even(n: Number): Boolean {
    if (n == 0) true else is_odd(n - 1);
}

function is_odd(n: Number): Boolean {
    if (n == 0) false else is_even(n - 1);
}
\end{lstlisting}

Las funciones builtin disponibles son: \texttt{print}, \texttt{sqrt},
\texttt{sin}, \texttt{cos}, \texttt{exp}, \texttt{log}, \texttt{rand} y
\texttt{range}. Además, \texttt{PI} y \texttt{E} se modelan como funciones
de aridad cero que actúan como constantes matemáticas.

\subsubsection{Orientación a objetos}

HULK implementa orientación a objetos con herencia simple y polimorfismo.
Los tipos se declaran con la palabra clave \texttt{type}, pueden recibir
parámetros de constructor e implementar métodos. Todos los atributos son
privados (accesibles solo dentro del tipo que los declara) y todos los métodos
son públicos y virtuales:

\begin{lstlisting}[caption={Tipos, herencia y polimorfismo en HULK}, label=lst:oop]
type Animal(n: String) {
    name: String = n;
    sound(): String { "..."; }
}

type Dog(n: String) inherits Animal(n) {
    sound(): String { "Woof"; }
}

type Cat(n: String) inherits Animal(n) {
    sound(): String { "Meow"; }
}

// Polimorfismo: variable de tipo Animal contiene un Dog
let a: Animal = new Dog("Rex") in print(a.sound()); // "Woof"
\end{lstlisting}

La palabra clave \texttt{base} permite delegar al método del padre inmediato:

\begin{lstlisting}[caption={Uso de base para delegación}, label=lst:base]
type FancyPrinter(prefix: String, suffix: String)
        inherits Printer(prefix) {
    sfx: String = suffix;
    format(msg: String): String {
        base(msg) @ self.sfx;   // delega a Printer.format y agrega sufijo
    }
}
\end{lstlisting}

\subsubsection{Protocolos}

Los protocolos en HULK definen interfaces estructurales: un tipo implementa un
protocolo si posee todos sus métodos con las signaturas correctas, sin necesidad
de declararlo explícitamente. Esto es \textit{tipado estructural} (similar a los
\textit{traits} de Go o las \textit{interfaces implícitas} de TypeScript):

\begin{lstlisting}[caption={Protocolos como interfaces estructurales}, label=lst:protocol]
protocol Printable {
    to_string(): String;
}

type Point(x: Number, y: Number) {
    x: Number = x;
    y: Number = y;
    to_string(): String { "point"; }
    // Point implementa Printable implícitamente
}

let p: Printable = new Point(1, 2) in
    print(p.to_string());
\end{lstlisting}

\subsubsection{Control de flujo}

HULK ofrece las construcciones de control clásicas, todas como expresiones:
\texttt{if}/\texttt{elif}/\texttt{else}, \texttt{while} y \texttt{for}.
El bucle \texttt{for} itera sobre cualquier expresión de tipo \texttt{Iterable<T>},
incluyendo rangos y vectores:

\begin{lstlisting}[caption={Control de flujo en HULK}, label=lst:control]
// if/elif/else como expresión
function grade(score: Number): String {
    if (score >= 90) "A"
    elif (score >= 80) "B"
    elif (score >= 70) "C"
    else "F";
}

// for sobre un rango
for (x in range(1, 6))
    print(x);
\end{lstlisting}

\subsection{El analizador léxico}

El lexer de HULK, implementado manualmente en el crate \texttt{hulk-lexer},
reconoce los siguientes grupos de tokens:

\begin{table}[h!]
\centering
\caption{Categorías de tokens del lexer de HULK}
\label{tab:tokens}
\begin{tabular}{@{}ll@{}}
\toprule
\textbf{Categoría} & \textbf{Tokens} \\
\midrule
Literales & \texttt{Number(f64)}, \texttt{StringLit(String)}, \texttt{True}, \texttt{False} \\
Aritmética & \texttt{Plus}, \texttt{Minus}, \texttt{Star}, \texttt{Slash}, \texttt{Caret}, \texttt{Percent} \\
Cadenas & \texttt{At} (\texttt{@}), \texttt{AtAt} (\texttt{@@}) \\
Comparación & \texttt{EqEq}, \texttt{Neq}, \texttt{Lt}, \texttt{Gt}, \texttt{Leq}, \texttt{Geq} \\
Booleanos & \texttt{And}, \texttt{Or}, \texttt{Not} \\
Asignación & \texttt{Assign} (\texttt{=}), \texttt{ColonEq} (\texttt{:=}) \\
Delimitadores & \texttt{LParen}, \texttt{RParen}, \texttt{LBrace}, \texttt{RBrace}, \texttt{LBracket}, \texttt{RBracket} \\
Separadores & \texttt{Semicolon}, \texttt{Comma}, \texttt{Colon}, \texttt{Dot} \\
Flechas & \texttt{Arrow} (\texttt{->}), \texttt{FatArrow} (\texttt{=>}) \\
Palabras clave & \texttt{let}, \texttt{in}, \texttt{if}, \texttt{elif}, \texttt{else}, \texttt{while}, \texttt{for}, \\
 & \texttt{function}, \texttt{type}, \texttt{inherits}, \texttt{new}, \texttt{self}, \\
 & \texttt{base}, \texttt{is}, \texttt{as}, \texttt{protocol}, \texttt{match}, \texttt{case} \\
\bottomrule
\end{tabular}
\end{table}

El lexer detecta y reporta los siguientes errores léxicos: cadenas sin terminar
(\texttt{UnterminatedString}), caracteres no reconocidos (\texttt{UnexpectedChar})
y secuencias de escape inválidas (\texttt{InvalidEscape}). Notably, la palabra
\texttt{base} es un \textit{keyword contextual}: el lexer la emite como
\texttt{TokenKind::Base}, pero el parser la acepta también como identificador
en posiciones de nombre de variable, atributo o parámetro —siguiendo el patrón
de palabras clave contextuales de C\# y SoulNG.

\newpage

% ============================================================
% SECCIÓN 4: EXTENSIONES PROPUESTAS
% ============================================================

\section{Extensiones propuestas}
\label{sec:extensiones}

Este proyecto introduce tres extensiones al lenguaje HULK original, todas
implementadas con modificaciones reales a la sintaxis, el AST y el análisis
semántico. Cada extensión está motivada por carencias concretas del lenguaje base
y se justifica mediante comparativas con lenguajes modernos.

% ─────────────────────────────────────────────────────────────
\subsection{Extensión 1: Expresiones \texttt{match}}
% ─────────────────────────────────────────────────────────────

\subsubsection{Nombre y propósito}

La \textbf{expresión \texttt{match}} introduce coincidencia de patrones
estructurada en HULK. El lenguaje base solo ofrece
\texttt{if}/\texttt{elif}/\texttt{else} para discriminar valores, lo que
resulta en código verboso cuando se necesita comparar un valor contra múltiples
casos. La expresión \texttt{match}, siendo una expresión (no una sentencia),
produce un valor que puede ser usado directamente en cualquier contexto.

\subsubsection{Sintaxis}

La gramática extendida para \texttt{match} es:

\begin{lstlisting}[language={},caption={Gramática BNF de match}, label=lst:match-bnf]
MatchExpr  ::= 'match' Expr '{' MatchCase+ '}'
MatchCase  ::= 'case' Pattern '=>' Expr ';'
Pattern    ::= '_'                          -- wildcard
             | Literal                     -- literal value
             | Identifier                  -- variable binding
             | TypeRef (',' Identifier)?   -- type pattern
\end{lstlisting}

Ejemplos de cada tipo de patrón:

\begin{lstlisting}[caption={Patrones literales}, label=lst:match-lit]
let x = 42 in
    match x {
        case 1  => print("uno");
        case 42 => print("cuarenta y dos");
        case _  => print("otro");
    };
\end{lstlisting}

\begin{lstlisting}[caption={Patrones de tipo con jerarquía de herencia},
                   label=lst:match-type]
type Animal { sound(): String { "..."; } }
type Dog inherits Animal { sound(): String { "Woof"; } }
type Cat inherits Animal { sound(): String { "Meow"; } }

let a: Animal = new Dog() in
    match a {
        case d: Dog => print(d.sound());   // narrows to Dog
        case c: Cat => print(c.sound());   // narrows to Cat
        case _      => print("unknown");
    };
\end{lstlisting}

\begin{lstlisting}[caption={match como expresión — el resultado se asigna},
                   label=lst:match-expr]
function describe(s: String): String {
    match s {
        case "hello" => "greeting";
        case "bye"   => "farewell";
        case _       => "unknown";
    }
}
print(describe("hello"));   // greeting
\end{lstlisting}

\subsubsection{Semántica estática}

El analizador semántico implementa la inferencia de \texttt{match} en la función
\texttt{infer\_match} del módulo \texttt{hulk-semantic}. Las reglas de tipado son:

\begin{enumerate}[leftmargin=*, label=\arabic*.]
    \item El escrutinio (expresión después de \texttt{match}) se infiere normalmente.
    \item Cada caso infiere su cuerpo en un entorno extendido según el patrón:
    \begin{itemize}
        \item \texttt{Wildcard} (\texttt{\_}): no añade bindings.
        \item \texttt{Literal}: verifica conformidad de tipos con el escrutinio.
        \item \texttt{Variable(name)}: enlaza \texttt{name} al tipo del escrutinio.
        \item \texttt{Type(T, binding)}: verifica que \texttt{T} sea un tipo
              conocido; si hay binding, lo enlaza con tipo \texttt{T} (narrowing).
    \end{itemize}
    \item El tipo del \texttt{match} completo es el \textit{ancestro común más
          bajo} (LCA) de los tipos de todos los cuerpos de caso.
    \item Si ningún caso es \textit{catch-all} (\texttt{\_} o variable), el
          compilador emite una advertencia \texttt{NonExhaustiveMatch}.
\end{enumerate}

\subsubsection{Semántica dinámica}

En tiempo de ejecución, la evaluación de \texttt{match x \{ case p => e \}}
procede así:

\begin{enumerate}[leftmargin=*, label=\arabic*.]
    \item Se evalúa \texttt{x} una sola vez y se guarda el resultado.
    \item Los casos se evalúan en orden de declaración.
    \item Para el primer patrón que coincida, se evalúa su cuerpo y se retorna
          el resultado.
    \item Si ningún patrón coincide, el runtime llama a
          \texttt{hulk\_rt\_match\_fail()} que termina el programa.
\end{enumerate}

El matching para patrones de tipo usa \texttt{hulk\_rt\_downcast\_check},
que recorre la cadena de vtables para verificar si el objeto es una instancia
del tipo esperado.

\subsubsection{Comparativa con otros lenguajes}

\begin{table}[h!]
\centering
\caption{Comparativa de match/switch en diferentes lenguajes}
\label{tab:match-compare}
\begin{tabularx}{\textwidth}{@{}lXX@{}}
\toprule
\textbf{Lenguaje} & \textbf{Construcción} & \textbf{Diferencias con HULK} \\
\midrule
Rust & \texttt{match} & Exhaustividad obligatoria; patrones destructurantes de structs/enums \\
Kotlin & \texttt{when} & Puede ser expresión o sentencia; soporte para rangos \\
Haskell & \texttt{case of} & Exhaustividad verificada; guardas con \texttt{|} \\
Python & \texttt{match} (3.10+) & Patrones estructurales; desestructuración de secuencias \\
Java & \texttt{switch} (14+) & Expresión switch con \texttt{->}; sin type narrowing automático \\
HULK & \texttt{match} & Siempre expresión; type narrowing; advertencia si no-exhaustivo \\
\bottomrule
\end{tabularx}
\end{table}

La decisión de emitir una \textit{advertencia} (no un error) para matches
no exhaustivos sigue el modelo de Kotlin, que también permite matches
no exhaustivos pero avisa al programador. Rust, en cambio, los rechaza
en compilación — una decisión más segura pero que dificulta el desarrollo
incremental.

% ─────────────────────────────────────────────────────────────
\subsection{Extensión 2: Tipos función y funciones de primera clase}
% ─────────────────────────────────────────────────────────────

\subsubsection{Nombre y propósito}

En HULK base, las funciones son entidades globales y no pueden almacenarse en
variables ni pasarse como argumentos. Esta extensión introduce el tipo
\textbf{función} \texttt{(T$_1$,...,T$_n$) -> R}, que permite tratar funciones
y métodos como valores de primera clase. Esto habilita la programación funcional
de alto orden: composición de funciones, paso de callbacks y referencias a métodos.

\subsubsection{Sintaxis}

El tipo función se escribe como:

\begin{lstlisting}[language={},caption={Sintaxis del tipo función},
                   label=lst:functype-bnf]
TypeRef ::= ... | '(' TypeList? ')' '->' TypeRef
\end{lstlisting}

Ejemplo de uso:

\begin{lstlisting}[caption={Función de alto orden con tipo función},
                   label=lst:functype-ex]
function apply(f: (Number) -> Number, x: Number): Number {
    f(x);
}

{
    if (apply(function (x: Number): Number -> x + 1, 5) == 6)
        print("ok")
    else
        print("fail");
}
\end{lstlisting}

Las referencias a métodos se obtienen accediendo al método sin llamarlo:

\begin{lstlisting}[caption={Referencia a método como valor},
                   label=lst:methodref]
type Counter(init: Number) {
    val: Number = init;
    inc(): Number { self.val := self.val + 1; self.val; }
}

let c = new Counter(0) in {
    let f: () -> Number = c.inc in {  // referencia al método
        print(f());   // llama c.inc()
        print(f());   // llama c.inc() de nuevo
    };
};
\end{lstlisting}

\subsubsection{Semántica estática}

El sistema de tipos extiende el enum \texttt{Type} con la variante:

\begin{lstlisting}[language=Rust, caption={Variante Type::Function en Rust},
                   label=lst:type-function]
pub enum Type {
    // ... tipos existentes ...
    /// Tipo función: (params) -> return_type.
    /// Usado para referencias a métodos y expresiones lambda.
    Function {
        params: Vec<Type>,
        return_type: Box<Type>,
    },
    // ...
}
\end{lstlisting}

La relación de subtipado para tipos función sigue la regla estándar de la
teoría de tipos: \textbf{contravarianza en parámetros, covarianza en retorno}:

\[
\frac{B_i \leq A_i \quad R_1 \leq R_2}
     {(A_1,...,A_n) \to R_1 \leq (B_1,...,B_n) \to R_2}
\]

Esta regla garantiza que una función más específica puede usarse donde se espera
una función más general.

Cuando el callee de una llamada tiene tipo \texttt{Type::Function}, el
analizador semántico extrae los tipos de parámetros y retorno, verifica la
aridad y los tipos de los argumentos, y produce un \texttt{CallExpr} tipado
con el tipo de retorno como anotación.

\subsubsection{Semántica dinámica}

Una referencia a método \texttt{obj.method} se representa en tiempo de ejecución
como un \textbf{thunk de dos palabras}: un puntero a la función y un puntero al
receptor (\texttt{self}). Esto permite llamar al método más tarde sin conocer el
tipo concreto del receptor.

\subsubsection{Comparativa con otros lenguajes}

\begin{table}[h!]
\centering
\caption{Funciones de primera clase en diferentes lenguajes}
\label{tab:functype-compare}
\begin{tabularx}{\textwidth}{@{}lXX@{}}
\toprule
\textbf{Lenguaje} & \textbf{Tipo función} & \textbf{Subtyping} \\
\midrule
Kotlin & \texttt{(T) -> R} & Contravariante/covariante; función vs. tipo \\
Rust & \texttt{fn(T) -> R} / \texttt{Fn(T) -> R} & Traits; closures capturan entorno \\
Haskell & \texttt{T -> R} & Curried; totalmente de primera clase \\
Java & \texttt{Function<T,R>} & Interfaces funcionales; lambdas desde Java 8 \\
TypeScript & \texttt{(x: T) => R} & Estructural; compatible por forma \\
HULK & \texttt{(T) -> R} & Contravariante/covariante; thunk para métodos \\
\bottomrule
\end{tabularx}
\end{table}

La decisión de usar un thunk de dos palabras para referencias a métodos sigue
el modelo de Kotlin y C\#, donde \texttt{obj::method} crea un delegate que
captura el receptor. En Rust, las referencias a métodos requieren closures
explícitas (\texttt{|x| obj.method(x)}), lo cual es más verboso.

% ─────────────────────────────────────────────────────────────
\subsection{Extensión 3: Vectores con comprensiones}
% ─────────────────────────────────────────────────────────────

\subsubsection{Nombre y propósito}

Esta extensión añade el tipo \textbf{\texttt{Vector<T>}} como colección
homogénea de tamaño fijo con soporte para:
\begin{itemize}
    \item Literales: \texttt{[1, 2, 3]}
    \item Indexación: \texttt{v[0]}
    \item Métodos: \texttt{v.size()}, \texttt{v.get(i)}, \texttt{v.set(i, x)}
    \item Comprensiones: \texttt{[expr | var in iterable]}
    \item Iteración con \texttt{for}
\end{itemize}

\subsubsection{Sintaxis}

\begin{lstlisting}[language={},caption={Gramática BNF de vectores},
                   label=lst:vec-bnf]
VectorExpr  ::= '[' VectorBody
VectorBody  ::= ']'
              | ExprNoTopOR VectorTail
VectorTail  ::= '|' id 'in' Expr ']'   -- comprehension
              | (',' Expr)* ']'          -- literal
TypeRef     ::= TypeRef '[]'            -- azucar: T[] = Vector<T>
              | TypeRef '*'             -- azucar: T* = Iterable<T>
\end{lstlisting}

\begin{lstlisting}[caption={Literal de vector e indexación},
                   label=lst:vec-lit]
let v = [1, 2, 3, 4, 5] in {
    print(v[0]);      // 1
    print(v.size());  // 5
};
\end{lstlisting}

\begin{lstlisting}[caption={Comprensión de vector al estilo funcional},
                   label=lst:vec-comp]
// Cuadrados de 1 a 5
let squares = [x * x | x in range(1, 6)] in
    for (s in squares)
        print(s);
// Salida: 1 4 9 16 25
\end{lstlisting}

Un detalle notable del parser es la resolución de la ambigüedad entre el
operador booleano \texttt{|} y el separador de comprensión. La gramática usa
una producción especial \texttt{ExprNoTopOR} para el lado izquierdo de la
comprensión, que no consume \texttt{|} en el nivel superior. Esto permite
que \texttt{[(x | y) | x in values]} sea válido porque el OR interno está
entre paréntesis.

\subsubsection{Semántica estática}

La inferencia de vectores implementada en \texttt{infer\_vector} sigue estas
reglas:

\begin{enumerate}[leftmargin=*, label=\arabic*.]
    \item \textbf{Literal}: el tipo del vector es \texttt{Vector<T>} donde
          \texttt{T} es el LCA de los tipos de todos los elementos.
    \item \textbf{Comprensión}: el tipo es \texttt{Vector<T>} donde \texttt{T}
          es el tipo de la expresión de elemento, después de inferir el tipo de
          la variable de iteración desde el iterable.
    \item \textbf{Indexación}: \texttt{v[i]} requiere que \texttt{v} sea
          \texttt{Vector<T>} e \texttt{i} sea \texttt{Number}; retorna \texttt{T}.
    \item \textbf{Covarianza}: \texttt{Vector<T>} $\leq$ \texttt{Vector<U>}
          si \texttt{T} $\leq$ \texttt{U}.
    \item \textbf{Iterabilidad}: \texttt{Vector<T>} implementa
          \texttt{Iterable<T>} automáticamente.
\end{enumerate}

\subsubsection{Comparativa con otros lenguajes}

\begin{table}[h!]
\centering
\caption{Colecciones y comprensiones en diferentes lenguajes}
\label{tab:vec-compare}
\begin{tabularx}{\textwidth}{@{}lXX@{}}
\toprule
\textbf{Lenguaje} & \textbf{Colección} & \textbf{Comprensión} \\
\midrule
Python & \texttt{list[T]} & \texttt{[x*x for x in range(5)]} \\
Haskell & \texttt{[T]} & \texttt{[x*x | x <- [1..5]]} \\
Kotlin & \texttt{List<T>} & \texttt{(1..5).map \{ it * it \}} \\
Rust & \texttt{Vec<T>} & \texttt{(1..6).map(|x| x*x).collect()} \\
Java & \texttt{List<T>} & Streams: \texttt{.stream().map(...).collect(...)} \\
HULK & \texttt{Vector<T>} & \texttt{[x*x | x in range(1,6)]} \\
\bottomrule
\end{tabularx}
\end{table}

La sintaxis de comprensión de HULK (\texttt{[expr | var in iter]}) es
deliberadamente similar a la de Haskell y Python, que son las más concisas
y legibles. La decisión de usar \texttt{|} como separador (en lugar de
\texttt{for} como en Python) sigue el modelo matemático de la notación
de conjuntos: $\{x^2 \mid x \in \mathbb{N}\}$.

\newpage

% ============================================================
% SECCION 5: ANALISIS DE DISENO Y ARQUITECTURA
% ============================================================

\section{An\'{a}lisis de dise\~{n}o y arquitectura}
\label{sec:arquitectura}

\subsection{Visi\'{o}n general del pipeline}

El compilador HULK implementado sigue la arquitectura cl\'{a}sica de compiladores
por fases, donde cada fase transforma una representaci\'{o}n del programa en otra
m\'{a}s elaborada. El pipeline completo es:

\begin{center}
\begin{tikzpicture}[
    node distance=0.4cm,
    box/.style={rectangle, draw=hulkblue, fill=hulkbg, rounded corners,
                minimum width=2cm, minimum height=0.7cm,
                font=\small\ttfamily, text centered},
    arrow/.style={->, thick, color=hulkblue}
]
\node[box] (src)  {fuente.hulk};
\node[box, right=of src]  (lex)  {hulk-lexer};
\node[box, right=of lex]  (par)  {hulk-parser};
\node[box, right=of par]  (sem)  {hulk-semantic};
\node[box, below=0.8cm of sem] (cg) {hulk-codegen};
\node[box, left=of cg]   (rt)   {hulk-rt};
\node[box, left=of rt]   (out)  {./output};
\draw[arrow] (src) -- (lex);
\draw[arrow] (lex) -- (par);
\draw[arrow] (par) -- (sem);
\draw[arrow] (sem) -- (cg);
\draw[arrow] (cg)  -- (rt);
\draw[arrow] (rt)  -- (out);
\end{tikzpicture}
\end{center}

El workspace de Rust est\'{a} organizado en siete crates con responsabilidades
claramente separadas, siguiendo el principio de responsabilidad \'{u}nica (SRP):

\begin{table}[h!]
\centering
\caption{Crates del workspace y sus responsabilidades}
\label{tab:crates}
\begin{tabularx}{\textwidth}{@{}lX@{}}
\toprule
\textbf{Crate} & \textbf{Responsabilidad} \\
\midrule
\texttt{hulk-ast}      & \'{A}rbol sint\'{a}ctico abstracto gen\'{e}rico parametrizado por el tipo de anotaci\'{o}n \\
\texttt{hulk-lexer}    & An\'{a}lisis l\'{e}xico: fuente a \texttt{Vec<Token>} \\
\texttt{hulk-parser}   & An\'{a}lisis sint\'{a}ctico LL(1): tokens a \texttt{Program<()>} \\
\texttt{hulk-semantic} & An\'{a}lisis sem\'{a}ntico en 5 pasadas: \texttt{Program<()>} a \texttt{VerifiedProgram} \\
\texttt{hulk-codegen}  & Generaci\'{o}n de c\'{o}digo LLVM IR, objeto y ejecutable \\
\texttt{hulk-rt}       & Biblioteca de runtime: GC, strings, vectores, builtins \\
\texttt{hulk-cli}      & Punto de entrada: orquesta el pipeline, gestiona exit codes \\
\bottomrule
\end{tabularx}
\end{table}

\subsection{El AST gen\'{e}rico parametrizado}

Una decisi\'{o}n de dise\~{n}o fundamental es que el AST (\texttt{hulk-ast}) est\'{a}
parametrizado por un tipo de anotaci\'{o}n \texttt{A}:

\begin{lstlisting}[language=Rust, caption={AST gen\'{e}rico en Rust}, label=lst:ast-generic]
pub struct Expr<A = ()> {
    pub kind: ExprKind<A>,
    pub anno: A,         // () antes del semantic; Type despues
    pub span: SourceSpan,
}
\end{lstlisting}

Esto permite que el mismo \'{a}rbol represente dos estados diferentes del programa:
\begin{itemize}
    \item \texttt{Program<()>} --- \'{a}rbol sin tipos, producido por el parser
    \item \texttt{Program<Type>} --- \'{a}rbol tipado, producido por el analizador sem\'{a}ntico
\end{itemize}

Esta t\'{e}cnica, conocida como \textit{typed syntax tree} o \textit{annotated AST},
es usada por compiladores profesionales como GHC (Haskell) y rustc. La ventaja
principal es que el compilador no puede mezclar accidentalmente \'{a}rboles tipados
y no tipados --- el sistema de tipos de Rust lo garantiza en tiempo de compilaci\'{o}n.

\subsection{El analizador l\'{e}xico}

El lexer implementado en \texttt{hulk-lexer} es un aut\'{o}mata finito determinista
(AFD) escrito manualmente. Lee la fuente car\'{a}cter a car\'{a}cter y emite tokens con
informaci\'{o}n de posici\'{o}n (\texttt{SourceSpan}).

Decisiones de dise\~{n}o:
\begin{itemize}
    \item \textbf{Lexer sin tabla}: cada estado es una rama del c\'{o}digo Rust,
          no una tabla de transiciones. Esto es m\'{a}s r\'{a}pido y m\'{a}s f\'{a}cil de
          depurar que los generadores de lexers como Flex.
    \item \textbf{Keywords como postproceso}: el lexer emite todos los
          identificadores como \texttt{Identifier}, luego un paso de
          postproceso convierte las palabras reservadas a sus tokens
          correspondientes.
    \item \textbf{\texttt{base} como keyword contextual}: resuelto en el parser,
          no en el lexer, siguiendo el patr\'{o}n C\#/SoulNG.
    \item \textbf{Errores recuperables}: el lexer reporta errores pero intenta
          continuar para reportar m\'{u}ltiples errores en una sola pasada.
\end{itemize}

\subsection{El analizador sint\'{a}ctico LL(1)}

El parser es un analizador descendente recursivo LL(1) implementado
manualmente. Cada no-terminal de la gram\'{a}tica corresponde a un m\'{e}todo de Rust:

\begin{lstlisting}[language=Rust, caption={Estructura del parser recursivo descendente},
                   label=lst:parser-struct]
impl Parser {
    fn parse_expression(&mut self) -> Result<Expr, ParseError> {
        self.parse_assignment()
    }
    fn parse_assignment(&mut self) -> Result<Expr, ParseError> {
        let lhs = self.parse_or()?;
        if self.peek_kind() == Some(&TokenKind::ColonEq) {
            self.advance();
            let rhs = self.parse_assignment()?;
            // construir AssignExpr...
        }
        Ok(lhs)
    }
}
\end{lstlisting}

Los puntos t\'{e}cnicos m\'{a}s importantes de la gram\'{a}tica son:

\begin{enumerate}[leftmargin=*, label=\arabic*.]
    \item \textbf{Eliminaci\'{o}n de recursi\'{o}n izquierda}: las reglas de precedencia
          ($E \to E \; op \; T \mid T$) se transforman a la forma iterativa
          ($E \to T \; E'$, $E' \to op \; T \; E' \mid \varepsilon$).
    \item \textbf{Resoluci\'{o}n de ambig\"{u}edad de vectores}: el operador \texttt{|}
          aparece tanto en booleanos como en comprensiones. Se resuelve con una
          producci\'{o}n especial \texttt{ExprNoTopOR}.
    \item \textbf{\texttt{base} contextual}: el helper \texttt{parse\_name()}
          acepta \texttt{TokenKind::Base} en posiciones de identificador, pero
          lo promueve a \texttt{ExprKind::BaseRef} solo como callee de llamada.
\end{enumerate}

\subsection{El analizador sem\'{a}ntico en cinco pasadas}

El an\'{a}lisis sem\'{a}ntico se organiza en cinco pasadas ordenadas, implementando el
patr\'{o}n \textit{two-pass binder} descrito por Immo Landwerth en el compilador Minsk:

\begin{table}[h!]
\centering
\caption{Las cinco pasadas del analizador sem\'{a}ntico}
\label{tab:passes}
\begin{tabularx}{\textwidth}{@{}clX@{}}
\toprule
\textbf{Pasada} & \textbf{Nombre} & \textbf{Responsabilidad} \\
\midrule
0 & \texttt{collect} & Registra todas las declaraciones globales en el
                       \texttt{TypeRegistry}. Permite referencias hacia adelante. \\
1 & \texttt{hierarchy} & Resuelve la jerarqu\'{i}a de herencia y conformidades de
                         protocolos. Detecta ciclos e incompatibilidades. \\
1.5 & \texttt{resolve\_constructor\_params} & Infiere tipos de par\'{a}metros de
                         constructor desde los inicializadores de atributos. \\
2 & \texttt{infer} & Inferencia de tipos bidireccional. Produce
                     \texttt{Program<Type>}. \\
3 & \texttt{check} & Verificaci\'{o}n final: conformidad de tipos, privacidad de
                     atributos, uso correcto de \texttt{self}. \\
\bottomrule
\end{tabularx}
\end{table}

La separaci\'{o}n entre \texttt{infer} (pasada 2) y \texttt{check} (pasada 3)
simplifica el c\'{o}digo: \texttt{infer} solo asigna tipos (puede producir
\texttt{Type::Unknown}), mientras que \texttt{check} opera sobre el \'{a}rbol ya
tipado y asume que todos los tipos est\'{a}n resueltos.

\subsection{El generador de c\'{o}digo LLVM}

El generador de c\'{o}digo usa LLVM 17 a trav\'{e}s de la biblioteca
\texttt{inkwell}. La funci\'{o}n de entrada p\'{u}blica es:

\begin{lstlisting}[language=Rust, caption={API p\'{u}blica de hulk-codegen},
                   label=lst:codegen-api]
pub fn compile(
    verified: &VerifiedProgram,
    opts: &CodegenOptions,
) -> Result<(), CodegenError>
\end{lstlisting}

\subsubsection{Modelo de objetos}

Cada tipo HULK se representa en memoria como una estructura con un
encabezado de objeto (\texttt{ObjHeader}) seguido de los atributos:

\begin{lstlisting}[language=C, caption={Layout de objetos en memoria},
                   label=lst:obj-layout]
struct ObjHeader {
    ref_count: i64,       // contador de referencias
    gc_mark:   i8,        // marca para mark-sweep
    next:      *ObjHeader, // lista del GC
    vtable:    *VTable,   // tabla de metodos virtuales
};
\end{lstlisting}

El \texttt{vtable} es un array de punteros a funciones en orden estable
(padre-antes-hijo, orden de declaraci\'{o}n). La llamada virtual consiste en:
cargar el puntero al vtable, indexar por el slot del m\'{e}todo, llamar indirectamente.

\subsubsection{Gesti\'{o}n de memoria}

El runtime implementa un modelo h\'{i}brido:
\begin{itemize}
    \item \textbf{Conteo de referencias}: cada asignaci\'{o}n incrementa el contador
          (\texttt{hulk\_rt\_retain}) y cada liberaci\'{o}n lo decrementa
          (\texttt{hulk\_rt\_release}). Cuando llega a cero, el objeto se libera.
    \item \textbf{Mark-sweep c\'{i}clico}: para manejar referencias circulares,
          el GC usa una \textit{shadow stack} de ra\'{i}ces. Cuando la memoria
          supera un umbral, \texttt{hulk\_rt\_gc\_collect()} libera objetos
          no alcanzables.
\end{itemize}

\subsection{Decisiones de dise\~{n}o no triviales}

\subsubsection{Rust como lenguaje de implementaci\'{o}n}

La elecci\'{o}n de Rust frente a C++ o Java est\'{a} justificada por:

\begin{enumerate}[leftmargin=*, label=\arabic*.]
    \item \textbf{Seguridad de memoria sin GC}: el sistema de tipos garantiza
          que no puede haber dangling pointers ni data races.
    \item \textbf{Enums algebraicos}: \texttt{Type}, \texttt{ExprKind} y
          \texttt{Pattern} se modelan como enums que fuerzan el manejo
          exhaustivo de todos los casos con \texttt{match}.
    \item \textbf{Ecosistema de tooling}: \texttt{cargo fmt},
          \texttt{cargo clippy} y \texttt{cargo test} proporcionan
          formateo, an\'{a}lisis est\'{a}tico y testing con zero configuraci\'{o}n.
\end{enumerate}

\subsubsection{Gram\'{a}tica LL(1) vs.\ LR}

Se eligi\'{o} un parser LL(1) recursivo descendente en lugar de un generador LR
por tres razones:

\begin{enumerate}[leftmargin=*, label=\arabic*.]
    \item \textbf{Mensajes de error}: los parsers LL(1) manuales producen
          mensajes m\'{a}s precisos porque el contexto de parsing es expl\'{i}cito.
    \item \textbf{Acciones sem\'{a}nticas directas}: cada m\'{e}todo construye el
          nodo del AST directamente, sin indirecci\'{o}n a trav\'{e}s de tablas.
    \item \textbf{Sin dependencias externas}: el proyecto se construye
          \'{i}ntegramente con \texttt{cargo}, sin necesitar flex/bison.
\end{enumerate}

\subsubsection{Separaci\'{o}n sem\'{a}ntica vs.\ codegen}

El dise\~{n}o separa expl\'{i}citamente el an\'{a}lisis sem\'{a}ntico del generador de
c\'{o}digo. La interfaz entre ambos es \texttt{VerifiedProgram}, que garantiza
que el codegen nunca recibe un programa con errores de tipo. Esta separaci\'{o}n
facilita el testing independiente de cada fase y permite apuntar a diferentes
backends (JVM, WASM, etc.) sin modificar el analizador sem\'{a}ntico.

\newpage

\end{document}